\begin{textsample}{NEG Dim 3 – 2022\_c5 – Score 2.00 – 2022\_c5/enade\_2022\_ri\_c2843d8b68.txt}  \label{ex:f3_neg_217}
CAPÍTULO 6 DISTRIBUIÇÃO DE CONCEITOS A finalidade deste capítulo é apresentar um panorama nacional sobre a distribuição do Conceito Enade 2022.
Para traçar esse panorama, recorre -se a tabelas e gráficos, que são acompanhados de comentários e análises que contribuem para a contextualização dos dados neles expostos.
Para visualizar o panorama do Enade/2022, em nível nacional, na primeira seção do capítulo, os dados sobre o número de cursos participantes são organizados por grande região.
Na segunda, o mapeamento do conceito do exame é feito por categoria administrativa das IES e modalidade de oferta dos cursos, de maneira cruzada às grandes regiões.
Na última seção, a lógica de apresentação dos dados é feita por organização acadêmica e, também, grande região de funcionamento dos cursos, no país como um todo. 6.1 PANORAMA NACIONAL DA DISTRIBUIÇÃO DOS CONCEITOS A lógica adotada na apresentação dos dados, nesta seção e nas duas seguintes, pressupõe o registro de duas informações de natureza reiterativa.
O primeiro diz respeito aos valores possíveis do Enade, que variam em níveis de 1 a 5, sendo que, à medida que esses valores aumentam, o curso é melhor avaliado e aqueles cursos com somente um ou sem nenhum concluinte participante não obtêm o Conceito Enade, ficam Sem Conceito ( SC ), como ressaltado na Apresentação desse relatório.
O segundo registro é que a diferença entre os cursos tabulados neste capítulo e no Capítulo 3 corresponde aos cursos que não foram avaliados, ou seja, àqueles, em princípio, sem estudantes concluintes que tivessem sido inscritos para a prova.
Na Tabela 6.1, são apresentadas a quantidade e a distribuição de cursos de Relações Internacionais participantes do Enade/2022, por faixa de conceito e grande região.
Observando -se os dados da Tabela 6.1, nota -se que, dos 131 cursos participantes, 45 ( 34,4 \% ) classificaram -se com conceito 3, o valor modal.
Esse foi também o conceito modal na maioria das regiões : Nordeste ( 40,0 \% ), Sul ( 55,2 \% ), Norte ( 66,7 \% ) e Centro-Oeste ( 41,7 \% ).
O conceito 2 foi o segundo mais frequente em nível nacional ( 32,1 \%, correspondendo a 42 cursos ), tendo sido o conceito modal na região Sudeste ( 39,2 \% ).
O conceito 4 foi o terceiro ( 18,3 \%, correspondendo a 24 cursos ).
Houve, ainda, 05 cursos ( 3,8 \% ) que receberam o conceito 5.
Dos 131 cursos de Relações Internacionais, 04 ( 3,1 \% ) ficaram Sem Conceito ( SC ).
A região Norte participou com 06 cursos.
Desses, 04 cursos receberam o conceito 3, o que equivale a 66,7 \% do total regional, o conceito modal na região, como já comentado.
Nenhum curso ficou Sem Conceito ( SC ), e dois cursos ( 33,3 \% ), obtiveram o conceito 2, nessa região.
Os conceitos 1, 4 e 5 não tiveram cursos atribuídos.
A região Nordeste participou com 10 cursos.
Nessa região, 04 cursos ( 40,0 \% em termos regionais ) obtiveram conceito 3, o conceito modal para a região, como já comentado.
Os conceitos 1 e 2 foram atribuídos, respectivamente, a 01 e a 02 cursos ( 10,0 \% e 20,0 \%, respectivamente ).
Ainda nessa região, 02 cursos ( 20,0 \% ) foram avaliados com o conceito 4, e um curso ( 10,0 \% ) foi avaliado com o conceito 5.
Nessa região, nenhum curso ficou Sem Conceito ( SC ).
Em um gráfico de distribuição cumulativa, em cada valor de abscissa ( eixo x ), o valor apresentado no eixo das ordenadas ( eixo y ) é o valor acumulado da variável ( Conceito Enade ).
Em linhas gerais, em um gráfico desta mesma tipologia, poligonais mais à esquerda representam, no caso dos dados apresentados nesta seção, as grandes regiões com a pior distribuição, e poligonais mais à direita, as grandes regiões com os melhores conceitos.
Para os 74 cursos participantes da região Sudeste, o conceito 2 foi o conceito modal, tendo sido obtido por 29 cursos ( 39,2 \% ).
O conceito 1 foi atribuído a oito cursos ( 10,8 \% ), e o conceito 3, a 16 cursos ( 21,6 \% ).
Receberam os conceitos 4 e 5, respectivamente, 17 cursos ( 23,0 \% ) e 1 curso ( 1,4 \% ).
Nessa região, três cursos ( 4,1 \% ) ficaram Sem Conceito ( SC ).
Os 29 cursos da região Sul corresponderam a 22,1 \% do total nacional.
Como já indicado, houve predominância do conceito 3 ( conceito modal ), atribuído a 55,2 \% dos cursos da região, ou seja, atribuído a 16 dos 29 cursos participantes na região Sul.
O conceito 1 foi atribuído a dois cursos ( 3,4 \% ) ; o conceito 2, a dois cursos ( 6,9 \% ) ; e os conceitos 4 e 5, a 03 ( 10,3 \% ) e a dois cursos ( 6,9 \% ), respectivamente.
Na região Sul, um curso ( 3,4 \% ) ficou Sem Conceito ( SC ).
Já dos 12 cursos participantes na região Centro-Oeste ( 9,7 \% do total nacional ), 05 receberam conceito 3, o conceito modal, como já destacado.
Quatro cursos ( 33,3 \% ) obtiveram o conceito 2, 02 cursos ( 16,7 \% ) obtiveram o conceito 4, e um curso ( 8,3 \% ), o conceito 5.
Nessa região, nenhum curso ficou Sem Conceito ( SC ).
O conceito 1 não teve cursos atribuídos.
Em um gráfico de distribuição cumulativa, em cada valor de abscissa ( eixo x ), o valor apresentado no eixo das ordenadas ( eixo y ) é o valor acumulado da variável ( Conceito Enade ).
Em linhas gerais, em um gráfico desta mesma tipologia, poligonais mais à esquerda representam, no caso dos dados apresentados nesta seção, as grandes regiões com a pior distribuição, e poligonais mais à direita, as grandes regiões com os melhores conceitos.
A partir destas diretrizes, no Gráfico 6.1, é apresentada a distribuição cumulativa do Conceito Enade, segundo a grande região.
Nesse gráfico, poligonais mais à esquerda representam Conceito Enade mais baixos, e poligonais mais à direita, Conceito Enade mais altos. 6.2 CONCEITOS POR CATEGORIA ADMINISTRATIVA, MODALIDADE DE OFERTA DOS CURSOS E GRANDE REGIÃO Na Tabela 6.2, é apresentada a distribuição dos cursos de Relações Internacionais participantes do Enade/2022, por categoria administrativa e por modalidade de oferta, de acordo com os conceitos por eles alcançados, segundo as grandes regiões brasileiras.
De acordo com as informações da Tabela 6.2, dos 131 cursos participantes, 28 eram ministrados em Instituições Públicas, e 103, em Instituições Privadas.
Quanto à modalidade de oferta, 12 cursos eram a Distância, e 119, Presenciais.
Ainda conforme as informações da mesma tabela, em termos nacionais, entre instituições públicas e privadas, observa -se que, dos 05 cursos avaliados com conceito 5, 03 foram oferecidos em IES públicas, e 02, em IES privadas.
Dos 28 cursos participantes de IES públicas, o conceito 3 foi o valor modal atribuído a 12 cursos.
Dos demais cursos de instituições públicas participantes, dois cursos foram avaliados com o conceito 2, 11 cursos foram avaliados com o conceito 4, e, 03 cursos, com o conceito 5.
Nessa categoria, nenhum curso ficou Sem Conceito ( SC ).
O conceito 1 não teve curso atribuído.
Na rede privada, o conceito modal foi o 2, atribuído a 40 cursos, o correspondente aos 119 cursos da categoria.
Dos demais cursos participantes, 11 receberam o conceito 1, e 33, o conceito 3.
O conceito 4 foi atribuído a 13 cursos, e o conceito 5, a 2 cursos.
Nessa categoria administrativa, 04 cursos ficaram Sem Conceito ( SC ).
Dos 119 cursos participantes da modalidade presencial, o conceito 3 foi o valor modal, atribuído a 42 cursos.
Dos demais cursos nessa modalidade de oferta, 08 cursos receberam o conceito 1 ; 38 cursos receberam o conceito 2 ; 24 cursos, o conceito 4 ; e 05 cursos, o conceito 5.
Nessa modalidade de oferta, dois cursos ficaram Sem Conceito ( SC ).
Dos 12 cursos participantes da modalidade a distância, o conceito 2 foi o valor modal, atribuído a 04 cursos.
Dos demais cursos nessa modalidade de oferta, 3 cursos receberam o conceito 1 ; outros 03 cursos receberam o conceito 3.
Nessa modalidade de oferta, dois cursos ficaram Sem Conceito ( SC ).
Na análise por região, observa -se que, na região Norte, tanto as instituições privadas quanto públicas, ambas tiveram os mesmos resultados, e participaram com 3 cursos, e nenhum curso ficou Sem Conceito ( SC ).
O conceito modal para as instituições privadas e públicas na região foi o 3, atribuído a dois cursos.
Um curso foi avaliado com o conceito 1 ; Nenhum curso ficou Sem Conceito ( SC ) ou recebeu os demais conceitos.
Dos cursos oferecidos na região Norte, 5 eram presenciais, e nenhum ficou Sem conceito ( SC ).
Dos 5 cursos oferecidos, dois obtiveram o conceito 2 ; 3 o conceito 3 ( modal ) ; Nenhum curso ficou Sem Conceito ( SC ) ou recebeu os demais conceitos.
Apenas um curso na modalidade a distância, tendo recebido o conceito modal 3.
Na região Nordeste, a rede privada concentrou 6 dos 10 cursos participantes.
Entre os cursos oferecidos por IES privadas no Nordeste, 3 cursos foram avaliados com conceito 3, o conceito modal.
Um curso para a região obteve os respectivos conceitos 1, 2 e 5.
Nenhum curso da região Sul oferecido por IES privadas ficou Sem Conceito ( SC ) ou recebeu o conceito 4.
Nessa categoria administrativa e grande região, 4 cursos ficaram Sem Conceito ( SC ).
As instituições públicas da região Nordeste participaram com 10 cursos.
Desses, dois cursos obtiveram conceito 4, e 4 cursos receberam conceito 3.
Nenhum curso da região Sul oferecido por IES púbica ficou Sem Conceito ( SC ) ou recebeu o conceito 1 e 5.
Dos cursos oferecidos na região Nordeste, sua totalidade foi para a modalidade presencial : 10.
Para esses, a moda foi o conceito 3, com 4 cursos.
Os demais cursos presenciais da região foram alocados nos conceitos 1 e 5, ambos com um curso, os conceitos 2 e 4.
Nenhum curso presencial da região ficou Sem Conceito ( SC ).
A modalidade a distância não teve custos para essa região.
Na região Sudeste, houve 74 cursos participantes.
Nessa combinação de categoria administrativa e grande região, o conceito modal foi o 2, atribuído a 29 cursos. 3 cursos ficaram Sem Conceito ( SC ).
Os demais foram avaliados com os seguintes conceitos : 1 ( oito cursos ), 3 ( 16 cursos ), 4 ( 17 cursos ) e 5 ( 1 curso ).
Dos 11 cursos oferecidos por instituições públicas na região Sudeste, o conceito modal foi o 4, atribuído a 8 cursos.
Três cursos receberam o conceito 3.
Nessa combinação de categoria administrativa e grande região, nenhum curso ficou Sem Conceito ( SC ), e nenhum recebeu o conceito 1, 2 e 5.
Também na região Sudeste, grande parcela dos cursos foi oferecida na modalidade presencial : 68 dos 74 cursos.
Para esses, a moda foi o conceito 2, com 27 cursos.
Os demais cursos presenciais da região foram alocados nos conceitos 1 ( seis cursos ), 3 ( 15 cursos ), 4 ( 17 cursos ) e 5 ( um curso ).
Ficaram Sem Conceito ( SC ) 2 cursos presenciais na região Sudeste.
Na região Sudeste, 6 cursos foram oferecidos na modalidade a distância, tendo os conceitos bimodal 1 e 2 ( dois cursos cada ) ; um curso com conceito 3.
Na região Sudeste, um curso a distância ficou Sem Conceito ( SC ).
As instituições privadas concentraram 22 dos 29 cursos participantes da região Sul.
Desses, 11 cursos obtiveram conceito 3, o conceito modal.
Nessa combinação de categoria administrativa e grande região, dois cursos receberam o conceito 1 ; 5 receberam o conceito 2 ; outros 3, o conceito 4 ; um curso ficou Sem Conceito ( SC ) e nenhum curso recebeu o conceito 5.
As instituições públicas na região Sul participaram com 7 cursos, aos quais foram atribuídos os conceitos 3 ( 5 cursos cada, modais ) e 5 ( 2 cursos ).
Nenhum curso da região Sul oferecido por IES públicas ficou Sem Conceito ( SC ) ou recebeu o conceito 1, 2 e 4. curso ), 2 ( 3 cursos ), 4 ( 3 cursos ), e 5 ( 2 cursos ).
Nenhum curso ficou Sem Conceito ( SC ).
Na modalidade a distância, foram oferecidos 5 cursos, sendo um curso atribuído ao conceito 1 ; 2 com conceito 2, o valor modal ; e 1 com conceito 3.
Um curso ficou Sem Conceito ( SC ) e nenhum com os conceitos 4 e 5 nesta região.
Na região Centro-Oeste, 9 dos 12 cursos participantes eram de IES.
Desses, 4 cursos a moda bimodal foram os conceitos 2 e 3.
Um curso recebeu o conceito 4.
Nenhum curso ficou Sem Conceito ( SC ) ou com os conceitos 1 e 5 nesta região.
Quanto aos 3 cursos oferecidos por instituições públicas na região Centro-Oeste, um curso foi atribuído, respetivamente, para os conceitos 3, 4 e 5.
Nenhum curso ficou Sem Conceito ( SC ) nessa região.
Também na região Centro-Oeste, a maioria dos cursos foi presencial : 12.
Para esses, a moda foi o conceito 3, com 5 cursos.
Os demais cursos presenciais da região foram alocados nos conceitos 2 ( quatro cursos ), 4 ( 2 cursos ) e 5 ( um curso ).
Nenhum curso presencial da região ficou Sem Conceito ( SC ).
A modalidade a distância não teve custos para essa região.
Sintetizando, no Gráfico 6.2, apresenta -se a distribuição do Conceito Enade, segundo a categoria administrativa da IES.
Nesse gráfico, como dito anteriormente, poligonais mais à esquerda representam Conceito Enade mais baixos, e poligonais mais à direita, Conceito Enade mais altos.
É importante observar que, para cada poligonal presente nesse gráfico, a porcentagem constante no eixo y se refere ao universo de cursos oferecidos por instituições públicas ( poligonal vermelha ) ou privadas ( poligonal azul ), respectivamente. 6.3 CONCEITOS POR ORGANIZAÇÃO ACADÊMICA E GRANDE REGIÃO Na Tabela 6.3, encontra -se a distribuição dos conceitos atribuídos aos cursos participantes do Enade/2022, na área de Relações Internacionais, por organização acadêmica, segundo as grandes regiões brasileiras.
Como mostram os dados da Tabela 6.3, dos 131 cursos de Relações Internacionais participantes, 82 eram oferecidos em universidades, 31 em centros universitários, 18 em faculdades e 0 em CEFET/IF.
De acordo com os dados apresentados na mesma tabela, de todos os 5 cursos avaliados com o conceito 5, 3 eram vinculados a universidades.
Os cursos vinculados a esse tipo de organização acadêmica tiveram o conceito 3 como conceito modal, atribuído a 31 cursos.
Os demais cursos ofertados pelas universidades receberam os conceitos 1 ( cinco cursos ), 2 ( vinte e um cursos ), 4 ( vinte e um cursos ) e conceito 5 ( três cursos, como já mencionado ).
Um curso ficou Sem Conceito ( SC ).
Dos cursos vinculados a centros universitários, o conceito modal foi o 2, atribuído a 14 cursos.
Dois cursos vinculados a esse tipo de organização acadêmica ficaram Sem Conceito ( SC ), e os outros cursos receberam os conceitos 1 ( 2 cursos ), 3 ( 10 cursos ), 4 ( 2 cursos ) e 5 ( um curso ).
Dos 18 cursos mantidos por faculdades, 7 receberam conceito 2, o conceito modal.
O segundo conceito mais frequentemente atribuído aos cursos oferecidos por essa categoria de IES foram os conceitos 1 e 3, recebidos por 4 cursos.
Na sequência, um curso obteve conceito 4 e um curso, o conceito 5.
Dos cursos oferecidos em faculdades, um ficou Sem Conceito ( SC ).
Nesta Organização acadêmica nenhum curso foi mantido pelo CEFET/IF.
Considerando -se, separadamente, as regiões brasileiras, verifica -se que, na região Norte, as universidades participaram com 5 cursos avaliados, aos quais foram atribuídos os seguintes conceitos : 2, a 1 curso ; e 3, o conceito modal, a 4 cursos.
Nenhum curso ficou Sem Conceito ( SC ) nessa combinação de categoria e região, nem com os conceitos 1, 4 e 5.
Nenhum curso foi atribuído a Centros universitários na região Norte, e nenhum curso oferecido em faculdades da região obteve os conceitos 1, 4 e 5.
Nessa região e Organização Acadêmica, nenhum curso foi oferecido no CEFET/IF.
Na região Nordeste, as universidades participaram com 6 dos 10 cursos da área de Relações Internacionais oferecidos.
O conceito trimodal foi 2, 3 e 4, atribuídos a 2 cursos cada.
Nenhum curso oferecido por universidades no Nordeste recebeu os conceitos 1 e 5 ou ficou Sem Conceito ( SC ).
Os centros universitários contaram com três cursos participantes na região Nordeste, 1 de cada recebeu os conceitos trimodais 1, 3 e 5.
Nenhum curso oferecido por centros universitários no Nordeste ficou Sem Conceito ( SC ) nem recebeu os conceitos 2 e 4.
As faculdades foram representadas por 1 curso na região Nordeste, que recebeu o conceito 3.
Nenhum curso ficou Sem Conceito ( SC ), e nenhum curso oferecido em faculdades da região obteve os conceitos 1, 2, 4 e 5.
Na região Sudeste, as universidades concentraram 43 dos 74 cursos de Relações Internacionais da região.
Dos cursos oferecidos em IES com esse tipo de organização acadêmica, nessa região, os conceitos bimodais foram 2 e 4, atribuídos a 14 cursos.
Os demais cursos receberam os conceitos 1 ( quatro cursos ) e 3 ( onze cursos ).
Nenhum curso ficou Sem Conceito ( SC ), e nenhum recebeu o conceito 5.
Os centros universitários participaram com 17 cursos na região Sudeste, dos quais 9 obtiveram o conceito modal 2.
O conceito 1 foi atribuído a um curso ; o conceito 3, a três cursos ; e o conceito 4, a dois cursos.
Dois cursos ficaram Sem Conceito ( SC ) nessa combinação de organização acadêmica e região, e nenhum recebeu o conceito 5.
Já as faculdades foram representadas por 14 cursos na região Sudeste, dos quais um ficou Sem Conceito ( SC ).
O conceito modal 2 foi atribuído a seis cursos.
Aos demais cursos foram atribuídos os conceitos 1 ( três cursos ), 3 ( dois cursos ), 4 ( um curso ) e 5 ( um curso ).
Dos 29 cursos da região Sul, 20 eram mantidos por universidades, para os quais o conceito modal foi 3, atribuído a onze cursos.
Os demais cursos receberam os conceitos 1 ( um curso ), 2 ( dois cursos ), 4 ( três cursos ) e 5 ( dois cursos ).
Em IES com esse tipo de organização acadêmica, um curso da região Sul ficou Sem Conceito ( SC ).
Dos 7 cursos participantes de centros universitários da região Sul, nenhum ficou Sem Conceito ( SC ), e quatro cursos receberam conceito 3, o valor modal.
O conceito restante ficou assim distribuído : três cursos ( conceito 2 ).
Quanto aos 2 cursos vinculados a faculdades na região Sul, somente um curso recebeu o conceito 1 e um curso recebeu o conceito 3.
Nenhum curso ficou Sem Conceito ( SC ), e nenhum recebeu os conceitos 1, 4 e 5.
Na região Centro-Oeste, 8 dos 12 cursos oferecidos eram mantidos por universidades.
Para esse tipo de Organização Acadêmica, o conceito modal foi 3, atribuído a três cursos ; e nenhum curso da região ficou Sem Conceito ( SC ).
Os demais cursos receberam os conceitos 2 ( dois cursos ), 4 ( dois cursos ) e 5 ( um curso ), nenhum recebeu o conceito 1.
Os centros universitários da região Centro-Oeste contaram com 4 cursos : dois cursos foram avaliados com o conceito 2 ; e dois cursos, com o conceito 3.
Nenhum curso ficou Sem Conceito ( SC ), e nenhum recebeu os conceitos 1, 4 e 5.
Nessa região e Organização Acadêmica, nenhum curso foi oferecido no CEFET/IF nem em faculdades.
De maneira semelhante à forma como se procedeu, em relação à seção anterior, encerra -se esta com o Gráfico 6.3, que apresenta a distribuição do Conceito Enade dos cursos de Relações Internacionais, segundo a organização acadêmica da IES.
No Gráfico 6.3, apresenta -se a distribuição do Conceito Enade, segundo a organização acadêmica da IES.
A leitura desse último gráfico é semelhante aos anteriores, seguindo -se a lógica de que poligonais deslocadas para a esquerda correspondem a Conceito Enade mais baixos ( 1, 2 e 3 ) e poligonais deslocadas para a direita correspondem a Conceito Enade mais altos.

% matched lemmas: 
\end{textsample}
